\subsection*{Чувствительность к пропущенным исходам и партиям выполнения}
Дополнительный анализ чувствительности выполнен после получения результатов и использует сохранённые записи. Принадлежность к партии восстановлена по идентификаторам назначений в манифесте продолжения: 460 назначений относятся к исходной партии, 540 --- к продолжению; сохранены 996 программ, четыре отправленных назначения не завершены. Таблица~\ref{tab:heldout-sensitivity-bounds} показывает границы парных разностей при всех возможных значениях неизвестных бинарных исходов. Совпадение нижней и верхней границ означает отсутствие неопределённости из-за пропусков в данном наборе, но не устраняет выборочную неопределённость. В табл.~\ref{tab:heldout-sensitivity-batch} приведены разности внутри партий на полных парах задач. Бутстрэп по задачам описывает выборочную неопределённость в этих наблюдаемых поднаборах; он не учитывает изменения сервиса во времени, распределение назначений между партиями и механизм пропусков. Число полных пар, чьи условия выполнены в разных партиях, равно 2, 1, 0 и 1 для SR--SN, MR--MN, MN--SN и MR--SR соответственно. Различия между партиями зависят также от состава задач и не устанавливают причинное влияние времени или провайдера. Поэтому результат исходной партии нельзя считать несмещённой оценкой того, что получилось бы при выполнении всех назначений в ней. Четыре заранее заданных теста Мак-Немара и поправка Холма сохранены.

\begin{table}[!htbp]\centering\small
\caption{Границы парных разностей при неизвестных исходах (процентные пункты).}\label{tab:heldout-sensitivity-bounds}
\begin{tabular*}{\textwidth}{@{\extracolsep{\fill}}lrr}\toprule
Контраст & 193 оцениваемые задачи & Все 200 задач \\ \midrule
SR$-$SN & $[-1.55,-0.52]$ & $[-5.00,+3.00]$ \\
MR$-$MN & $[+1.04,+1.04]$ & $[-2.50,+4.50]$ \\
MN$-$SN & $[-4.66,-3.63]$ & $[-8.00,+0.00]$ \\
MR$-$SR & $[-2.07,-2.07]$ & $[-5.50,+1.50]$ \\
\bottomrule\end{tabular*}\par\smallskip\footnotesize Каждый неизвестный исход может принимать значение 0 или 1. Показаны границы возможных значений, а не доверительные интервалы.\end{table}

\begin{table}[!htbp]\centering\scriptsize
\caption{Чувствительность качества внутри партий (процентные пункты).}\label{tab:heldout-sensitivity-batch}
\begin{tabular*}{\textwidth}{@{\extracolsep{\fill}}llrrr}\toprule
Партия & Контраст & $N$ & Разность & 95\%-интервал бутстрэпа \\ \midrule
Исходная & SR$-$SN & 87 & -2.30 & $[-8.05,+3.45]$ \\
Исходная & MR$-$MN & 89 & +3.37 & $[-3.37,+11.24]$ \\
Исходная & MN$-$SN & 87 & -2.30 & $[-9.20,+4.60]$ \\
Исходная & MR$-$SR & 90 & +4.44 & $[-2.22,+11.11]$ \\
Продолжение & SR$-$SN & 102 & +0.00 & $[-4.93,+4.90]$ \\
Продолжение & MR$-$MN & 103 & -0.97 & $[-5.83,+3.88]$ \\
Продолжение & MN$-$SN & 104 & -6.73 & $[-13.46,+0.00]$ \\
Продолжение & MR$-$SR & 102 & -7.84 & $[-13.73,-2.94]$ \\
\bottomrule\end{tabular*}\par\smallskip\footnotesize Интервалы описывают выборочную неопределённость по задачам среди наблюдаемых полных пар; они не учитывают изменения сервиса во времени и механизм пропусков. Превышение времени нативной проверки считается наблюдаемым неуспехом согласно исходному протоколу.\end{table}
